\documentclass{report}
\usepackage{amsmath,amssymb}
\setlength{\parindent}{0mm}
\setlength{\parskip}{1em}
\begin{document}
\begin{center}
\rule{6in}{1pt} \
{\large James Adler \\
{\bf A Deflation Technique for Detecting Multiple Liquid Crystal Equilibrium States}}

Tufts University \\ 503 Boston Ave \\ Medford \\ MA 02155
\\
{\tt james.adler@tufts.edu}\\
David B.  Emerson\\
Patrick E.  Farrell\\
Scott P.  MacLachlan\end{center}

Multiple equilibrium states arise in many physical systems, including
various types of liquid crystal structures. Having the ability to
reliably compute such states enables more accurate physical analysis and
understanding of experimental behavior. In this talk, we consider
adapting and extending a deflation technique for the computation of
multiple distinct solutions arising in the context of modeling
equilibrium configurations of nematic and cholesteric liquid crystals.
The deflation method is applied as part of an overall free-energy
variational approach and is modified to fit the framework of optimization
of a functional with pointwise constraints. It is shown that multigrid
methods designed for the undeflated systems may be applied to efficiently
solve the linear systems arising in the application of deflation. For the
numerical algorithm, the deflation approach is interwoven with nested
iteration, creating a dynamic and efficient method that further enables
the discovery of distinct solutions. Finally, we present numerical
simulations demonstrating the efficacy and accuracy of the algorithm in
detecting important physical phenomena, including bifurcations and
disclination behaviors.


\end{document}
